% docs/slides/preamble.tex -- 五份 deck 共享
\documentclass[aspectratio=169,11pt]{beamer}

% ==== 中文与字体 ====
\usepackage[UTF8,fontset=none]{ctex}
\usepackage{fontspec}

\IfFontExistsTF{PingFang SC}{%
  \setCJKmainfont{PingFang SC}[BoldFont=PingFang SC,ItalicFont=PingFang SC,AutoFakeSlant=0.15]
  \setCJKsansfont{PingFang SC}[ItalicFont=PingFang SC,AutoFakeSlant=0.15]
  \setCJKmonofont{PingFang SC}
}{%
  \IfFontExistsTF{Source Han Sans SC}{%
    \setCJKmainfont{Source Han Sans SC}
    \setCJKsansfont{Source Han Sans SC}
  }{%
    \IfFontExistsTF{Noto Sans CJK SC}{%
      \setCJKmainfont{Noto Sans CJK SC}
      \setCJKsansfont{Noto Sans CJK SC}
    }{%
      \setCJKmainfont{Heiti SC}
      \setCJKsansfont{Heiti SC}
    }
  }
}

\IfFontExistsTF{Helvetica Neue}{\setmainfont{Helvetica Neue}}{%
  \IfFontExistsTF{Inter}{\setmainfont{Inter}}{}}
\IfFontExistsTF{JetBrains Mono}{\setmonofont{JetBrains Mono}[Scale=0.85]}%
  {\setmonofont{Menlo}[Scale=0.85]}

% ==== 颜色 ====
\usepackage{xcolor}
\definecolor{cMain}{HTML}{1F4E79}
\definecolor{cAccent}{HTML}{E07A1F}
\definecolor{cGray}{HTML}{595959}
\definecolor{cMute}{HTML}{F4F4F4}
\definecolor{cOK}{HTML}{2E7D32}
\definecolor{cBad}{HTML}{B23A3A}
\definecolor{cWarn}{HTML}{C8A300}

% ==== 主题 ====
\usepackage{tcolorbox}
\usetheme{metropolis}
\metroset{progressbar=frametitle,numbering=fraction,block=fill}
\setbeamercolor{frametitle}{bg=cMain,fg=white}
\setbeamercolor{progress bar}{fg=cAccent,bg=cMute}
\setbeamercolor{alerted text}{fg=cAccent}
\setbeamercolor{title}{fg=cMain}

% metropolis 默认尝试 Fira Sans，未安装时 fallback 到 Latin Modern Sans。
% 这里在主题加载后显式覆盖，强制使用 macOS 上一定存在的 Helvetica Neue，
% 跟中文 PingFang 风格一致；找不到则按顺序尝试 Inter / Arial。
\IfFontExistsTF{Helvetica Neue}{%
  \setsansfont{Helvetica Neue}%
}{%
  \IfFontExistsTF{Inter}{\setsansfont{Inter}}{%
    \IfFontExistsTF{Arial}{\setsansfont{Arial}}{}}%
}

% ==== TikZ ====
\usepackage{tikz}
\usetikzlibrary{positioning,arrows.meta,fit,shapes.geometric,calc,%
  decorations.pathmorphing,backgrounds}

\tikzset{
  node-box/.style={draw=cMain,thick,rounded corners=2pt,
    minimum height=8mm,minimum width=18mm,align=center,font=\small},
  node-state/.style={draw=cMain,thick,circle,minimum size=14mm,
    align=center,font=\small},
  node-ok/.style={draw=cOK,thick,fill=cOK!10},
  node-warn/.style={draw=cWarn,thick,fill=cWarn!15},
  node-bad/.style={draw=cBad,thick,fill=cBad!10},
  % 色盲友好的状态机三态：蓝 (正常/Closed) / 橙 (告警/Open) / 灰 (中间/HalfOpen)
  % 避免 red-green 二元区分，改用 hue + lightness 双通道。
  node-state-normal/.style={draw=cMain,thick,fill=cMain!12},
  node-state-alert/.style={draw=cAccent,thick,fill=cAccent!18},
  node-state-transition/.style={draw=cGray,thick,fill=cGray!18},
  arrow/.style={-{Stealth[length=2mm]},thick,draw=cMain},
  arrow-async/.style={-{Stealth[length=2mm]},thick,dashed,draw=cGray},
  arrow-bad/.style={-{Stealth[length=2mm]},thick,draw=cBad},
  arrow-alert/.style={-{Stealth[length=2mm]},thick,draw=cAccent},
  lifeline/.style={dashed,thick,draw=cGray},
}

% ==== 代码 ====
\usepackage{listings}
\lstdefinelanguage{Go}{
  morekeywords={break,case,chan,const,continue,default,defer,else,
    fallthrough,for,func,go,goto,if,import,interface,map,package,range,
    return,select,struct,switch,type,var,nil,true,false,iota},
  morekeywords=[2]{string,int,int64,uint,uint64,bool,byte,error,context},
  sensitive=true,
  morecomment=[l]//,morecomment=[s]{/*}{*/},
  morestring=[b]",morestring=[b]`,
}
\lstdefinelanguage{LuaRedis}{
  morekeywords={and,break,do,else,elseif,end,false,for,function,if,in,
    local,nil,not,or,repeat,return,then,true,until,while},
  morekeywords=[2]{redis,call,KEYS,ARGV,tonumber,HGET,HSET,GET,SET,
    DECR,DECRBY,INCR,INCRBY,EXPIRE,PEXPIRE,ZADD,ZCARD,ZREMRANGEBYSCORE,
    EVAL,EVALSHA},
  sensitive=true,
  morecomment=[l]{--},
  morestring=[b]",morestring=[b]',
}

\lstdefinestyle{base}{
  basicstyle=\ttfamily\scriptsize,
  numbers=left,numberstyle=\tiny\color{cGray},
  keywordstyle=\color{cMain}\bfseries,
  keywordstyle=[2]\color{cAccent},
  stringstyle=\color{cOK},
  commentstyle=\color{cGray}\itshape,
  backgroundcolor=\color{cMute},
  frame=single,framerule=0pt,
  showstringspaces=false,
  breaklines=true,
  tabsize=2,
  columns=fullflexible,
}
\lstset{style=base}

% ==== 三线表与附录页码 ====
\usepackage{booktabs}
\usepackage{appendixnumberbeamer}

% ==== 页脚来源标注 ====
\newcommand{\srcnote}[1]{{\color{cGray}\scriptsize\itshape #1}}

% 关键论点框
\newcommand{\keypoint}[1]{%
  \begin{tcolorbox}[colback=cMute,colframe=cMain,boxrule=0.5pt,arc=2pt]%
  #1\end{tcolorbox}}


\title{商家在 gomall 里是怎么活的}
\subtitle{BossID 字段在、商家自助 API 缺：从"商家用 admin 后台"到"商家用商家后台"}
\author{gomall · 业务侧 deck}
\date{}

\begin{document}

\maketitle

\begin{frame}{今日大纲}
\tableofcontents
\end{frame}

% ============================================================
\section{业务困局：商家在 gomall 里没有"自己的后台"}
% ============================================================

\begin{frame}{一句话：BossID 字段早就在了，商家接口一个都没有}
\begin{itemize}
  \item gomall 五张核心表里，\textbf{BossID 字段都已经埋好}：order、product、cart、favorite、skill\_product。
  \item 但全站没有任何一个 \texttt{/api/v1/merchant/*} 路由 —— 商家想看自己店铺数据\textbf{进不去}。
  \item 现状导致三件事：
  \begin{itemize}
    \item 商家想看"我今天卖了多少" $\to$ 没接口，要么找客服查、要么直接看 admin 后台（越权）。
    \item 商家想知道"我哪个 SKU 快没库存" $\to$ 没告警，等用户下不了单才发现。
    \item 商家上架商品调 \texttt{product/create} —— 这接口在 authed group 下，\textbf{任何登录用户都能上架}。
  \end{itemize}
\end{itemize}
\srcnote{routes/router.go:60--62 product/create 三件套都在 authed 下，没 RequireRole；:135--141 admin sub-group 只有 3 个 handler}
\end{frame}

\begin{frame}{现状对业务的真实伤害}
\begin{itemize}
  \item \textbf{招商困难}：BD 谈一个新商家，洽谈阶段只能开 admin 账号给他 —— 商家立刻看到\textbf{所有店铺}数据，谈判桌上当场翻车。
  \item \textbf{客服压力}：商家把客服当后台用 —— "帮我查下我店铺今日订单"、"帮我看下 SKU 12345 的库存"。客服 1 个人对 N 个商家，挤兑出的工单全是低价值查询。
  \item \textbf{权限边界没法解释}：\texttt{product/delete} 已经有 \texttt{WHERE id=? AND boss\_id=?}，但\textbf{没任何中间件保证"调用者一定是 boss\_id 这个人"} —— 数据查询正确，权限模型却没成形。
\end{itemize}
\srcnote{repository/db/dao/product.go:67--72 DeleteProduct 已带 boss\_id 隔离；但 router.go:62 没有 RequireMerchant}
\end{frame}

% ============================================================
\section{商家用户旅程：哪一站有路、哪一站没路}
% ============================================================

\begin{frame}{商家七步旅程 + gomall 现状}
\begin{tikzpicture}[node distance=1mm and 1mm,every node/.append style={align=center,font=\tiny,text width=13mm}]
  \node[node-box,node-state-normal] (s1) {1.注册};
  \node[node-box,node-state-alert,right=of s1] (s2) {2.资质\\审核};
  \node[node-box,node-state-alert,right=of s2] (s3) {3.第一次\\上架};
  \node[node-box,node-state-normal,right=of s3] (s4) {4.第一单\\成交};
  \node[node-box,node-state-alert,right=of s4] (s5) {5.售后\\退款};
  \node[node-box,node-state-alert,right=of s5] (s6) {6.结算\\提现};
  \node[node-box,node-state-alert,right=of s6] (s7) {7.数据\\复盘};
  \draw[arrow] (s1) -- (s2); \draw[arrow] (s2) -- (s3); \draw[arrow] (s3) -- (s4);
  \draw[arrow] (s4) -- (s5); \draw[arrow] (s5) -- (s6); \draw[arrow] (s6) -- (s7);
\end{tikzpicture}
\vspace{1mm}
\scriptsize
\begin{tabular}{@{}clll@{}}
\toprule
站 & 业务诉求 & gomall 现状 & 命中表/接口 \\
\midrule
1 & 注册        & \checkmark 复用 user/register & user 表 Role='user' \\
2 & 资质审核    & 缺            & 无 merchant\_profile 表 \\
3 & 上架商品    & 部分          & product/create 任谁都能调 \\
4 & 第一单成交  & \checkmark 但商家看不到 & order 表已有 boss\_id \\
5 & 售后退款    & 缺            & 无商家侧退款审批 \\
6 & 结算提现    & 缺            & 无 settlement / withdraw \\
7 & 数据复盘    & 缺            & 无销量 / 转化看板 \\
\bottomrule
\end{tabular}
\vspace{1mm}
\small \textbf{"部分"的代价}：\texttt{product/create} 无角色校验，任何注册用户调一次自己就成了"店主"。gomall 业务定位是 B2C 自营+第三方品牌\textbf{不是闲鱼}；上架前必须资质审核。
\srcnote{repository/db/model/order.go:12 BossID；user.go:23 Role 只有 user/admin 两态；router.go:60 product/create}
\end{frame}

% ============================================================
\section{三层鉴权边界：把 router 翻译成商家视角}
% ============================================================

\begin{frame}{三层墙：商家应该站在中间那层}
\begin{tikzpicture}[node distance=3mm and 3mm,every node/.append style={align=center,font=\scriptsize,text width=26mm}]
  \node[node-box,node-state-normal] (user)  {C 端用户\\\texttt{/api/v1/*}\\下单/领券/支付};
  \node[node-box,node-state-transition,right=of user] (merchant) {\textbf{商家（缺）}\\\texttt{/api/v1/merchant/*}\\自店订单/商品/库存};
  \node[node-box,node-state-alert,right=of merchant] (admin) {admin\\\texttt{/api/v1/admin/*}\\promote/backfill};
  \draw[arrow]      (user)     -- node[above,font=\tiny]{merchant 角色} (merchant);
  \draw[arrow-alert](merchant) -- node[above,font=\tiny]{提权审核} (admin);
\end{tikzpicture}
\vspace{1mm}
\scriptsize
\begin{tabular}{@{}lllll@{}}
\toprule
角色 & 能看 & 能改 & 路径 & 现状 \\
\midrule
C 端用户  & 自己的订单         & 自己资料   & \texttt{/api/v1/*}          & \checkmark \\
商家     & 自店订单/商品/库存 & 自店范围   & \texttt{/api/v1/merchant/*} & \textbf{缺} \\
admin   & 全部              & 全部      & \texttt{/api/v1/admin/*}    & 部分 \\
\bottomrule
\end{tabular}
\vspace{1mm}
\small \textbf{三层各自的最坏情况}：C 端越权 $\to$ 看到别人订单（泄露电话/地址）；商家越权 $\to$ 看到其他商家销售（竞业/抄品）；admin 越权 $\to$ 一切（必须审计，见 deck 14）。
\srcnote{routes/router.go:46 authed；:135--141 admin；middleware/rbac.go:48--86 RequireRole；model/user.go:30--31 Role 常量}
\end{frame}

% ============================================================
\section{数据基础：BossID 字段在哪}
% ============================================================

\begin{frame}{BossID 散落在 5 张表 + 现有 dao 用法}
\scriptsize
\begin{tabular}{@{}lll@{}}
\toprule
位置 & 字段或语句 & 业务含义 \\
\midrule
model/order.go:12       & BossID uint                    & 这单卖家是谁 \\
model/product.go:23     & BossID uint                    & 这商品归谁 \\
model/cart.go:10        & BossID uint                    & 加购的店铺归属 \\
model/favorite.go:11    & BossID uint                    & 收藏背后的店主 \\
model/skill\_product.go:6& BossId uint                   & 秒杀活动发起商家 \\
dao/product.go:69       & DELETE WHERE id AND boss\_id   & 商家删自己商品 \\
dao/cart.go:61          & SELECT WHERE user/product/boss & 加购"店内唯一性" \\
dao/order.go:135        & o.boss\_id AS boss\_id (JOIN)  & 订单详情带卖家 \\
dao/favories.go:33      & JOIN user ON u.id=f.boss\_id   & 关注的店铺 \\
\bottomrule
\end{tabular}
\vspace{1mm}
\small \textbf{所有 boss\_id 用法都是"C 端用户视角"} —— 拿来展示"这商品/订单/购物车的卖家是谁"，\textbf{从来没有"以 boss\_id 为主语"的查询}。
\keypoint{有字段 $\ne$ 有能力。商家自助接口的本质是把 boss\_id \textbf{从外键升级为查询主体}。}
\end{frame}

% ============================================================
\section{商家后台 MVP：第一版接口设计}
% ============================================================

\begin{frame}[fragile]{接口设计草案：沿用 admin group 写法}
\begin{lstlisting}[language=Go,firstnumber=135,
  caption={\srcnote{对照 routes/router.go:135--141 admin 子组，新增 merchant 子组}}]
// 商家子组：先 authed，再 RequireRole("merchant")
merchant := authed.Group("/merchant")
merchant.Use(middleware.RequireRole("merchant"))
{
    merchant.GET ("orders",       MerchantListOrders)      // WHERE boss_id = ctx.userID
    merchant.GET ("orders/stats", MerchantOrderStats)      // 日/周/月销量
    merchant.GET ("products",     MerchantListProducts)
    merchant.POST("products",     MerchantCreateProduct)   // 替换裸 product/create
    merchant.GET ("inventory",    MerchantInventoryStatus) // 低库存告警
    merchant.POST("coupon/batch", MerchantCreateCoupon)
    merchant.GET ("settlement",   MerchantSettlement)      // 结算/提现
}
\end{lstlisting}
\small \textbf{照搬 admin 写法}：\texttt{authed.Group} + \texttt{Use(RequireRole)}，复用 30s 角色缓存；\texttt{RequireRole} 已支持变参，无需改实现。
\srcnote{middleware/rbac.go:48--86 RequireRole(allowed...string)}
\end{frame}

\begin{frame}{MVP 七接口的业务优先级}
\begin{tabular}{@{}lll@{}}
\toprule
接口 & 业务紧迫度 & 实现工作量 \\
\midrule
GET orders            & 高（替代客服查单） & 小（已有 order dao） \\
GET orders/stats      & 高（招商对外）     & 中（新增聚合 SQL） \\
GET inventory         & 高（防超卖恶化）   & 小（复用库存表） \\
POST coupon/batch     & 中（已有，需限店） & 小（加 WHERE boss\_id） \\
GET products          & 中                & 小 \\
POST products         & 中（替换裸接口）   & 小（迁移 + 加 RBAC） \\
GET settlement        & 低（结算还没做）   & 大（涉及对账 / 第三方） \\
\bottomrule
\end{tabular}
\vspace{2mm}
\small \textbf{业务侧建议顺序}：先做\textbf{读接口}（让商家\textbf{看得到}），再做\textbf{告警接口}（让商家\textbf{被通知到}），最后做\textbf{结算接口}（让商家\textbf{拿得到钱}）。三步对应"信任三段论"：能看 $\to$ 能管 $\to$ 能赚。
\end{frame}

% ============================================================
\section{商家看板：信息架构四象限}
% ============================================================

\begin{frame}{商家看板四象限 + SQL 草图}
\begin{columns}[T,onlytextwidth]
\column{0.42\textwidth}
\begin{tikzpicture}[node distance=2mm and 2mm,every node/.append style={align=center,font=\tiny,text width=22mm}]
  \node[node-box,node-state-normal] (ne)  {\textbf{今日营收}\\GMV / 客单价};
  \node[node-box,node-state-alert,right=of ne] (nw) {\textbf{待发货}\\type=2 未发};
  \node[node-box,node-state-transition,below=of ne] (se) {\textbf{库存告警}\\num<safety};
  \node[node-box,node-state-alert,below=of nw] (sw) {\textbf{客诉}\\退款率/待回};
\end{tikzpicture}
\column{0.55\textwidth}
\scriptsize
\begin{tabular}{@{}p{18mm}p{42mm}@{}}
\toprule
象限 & 查询 \\
\midrule
今日营收 & \texttt{SELECT COUNT,SUM(money) FROM order WHERE boss\_id=? AND DATE=CURDATE AND type=2} \\
待发货   & \texttt{WHERE boss\_id=? AND type=2 ORDER BY created\_at} \\
库存告警 & \texttt{WHERE boss\_id=? AND num < safety\_stock} \\
客诉     & \texttt{aftersale\_ticket WHERE boss\_id=?}（待建）\\
\bottomrule
\end{tabular}
\end{columns}
\vspace{2mm}
\small \textbf{已有索引}：\texttt{order.user\_id} 走索引（C 端用），\texttt{boss\_id 没建索引}。商家维度查询前\textbf{先 ALTER TABLE order ADD INDEX idx\_boss\_date(boss\_id, created\_at)}，否则像 deck 13 那样翻全表。
\srcnote{repository/db/model/order.go:16 Type 1未支付/2已支付；product.go:22 Num 即可用库存；dao/order.go:50 现有 user\_id 索引参考}
\end{frame}

% ============================================================
\section{商家告警的业务规则}
% ============================================================

\begin{frame}{三条核心告警规则 + 阈值设计}
\scriptsize
\begin{tabular}{@{}llll@{}}
\toprule
规则 & 触发 & 通道 & 接收方 \\
\midrule
库存阈值 & num < safety\_stock        & outbox \texttt{inventory.low}  & 商家钉钉 \\
退款异常 & 单日退款率 > 10\%          & outbox \texttt{aftersale.spike}& 平台+商家 \\
异常订单 & 单订单金额 > 均值 5$\times$ & outbox \texttt{order.suspect}  & 风控审查 \\
\bottomrule
\end{tabular}
\vspace{1mm}
\small
\begin{itemize}
  \item 三条规则共用机制：\textbf{业务事务里写 outbox} $\to$ \texttt{publisher.loop} 异步投递。同步发短信 = 网关挂一次下单全卡死。
  \item \textbf{阈值不是 0}：库存等于 0 告警 = 已卖断货 = 用户买不到。要用 \texttt{safety\_stock}（按品类 = 7 天销量）。
  \item \textbf{safety\_stock 谁定}：商家自填大部分填 0；平台算准但要 7 天历史。\textbf{折中}：上架前 7 天按品类均值兜底，7 天后切实际销量。
  \item \textbf{频率收敛}：同 SKU 24h 只发一次，防商家被噪音淹没。
\end{itemize}
\srcnote{service/outbox/publisher.go:65--86 drainBatch；model/product.go:22 Num；docs/slides/04-outbox-saga.tex 配套}
\end{frame}

% ============================================================
\section{商家结算与提现：资金路径}
% ============================================================

\begin{frame}{T+1 结算 + 隔离账户}
\begin{tikzpicture}[node distance=3mm and 8mm]
  \node[node-state,node-state-normal] (s1) {D 日\\签收};
  \node[node-state,node-state-transition,right=of s1] (s2) {D+1 0:00\\结算};
  \node[node-state,node-state-transition,right=of s2] (s3) {D+1\\提现};
  \node[node-state,node-state-alert,right=of s3] (s4) {D+2\\出账};
  \draw[arrow]       (s1) -- node[above,font=\tiny]{冻结期} (s2);
  \draw[arrow]       (s2) -- node[above,font=\tiny]{settlement} (s3);
  \draw[arrow-alert] (s3) -- node[above,font=\tiny]{银行} (s4);
\end{tikzpicture}
\vspace{1mm}
\scriptsize
\begin{tabular}{@{}lll@{}}
\toprule
账户       & 含义              & 状态 \\
\midrule
冻结户 A   & 已支付未发货      & 不可结算 \\
冻结户 B   & 已发货未签收      & 不可结算 \\
可结算户   & 已签收 + 过退货期 & 可发起提现 \\
\bottomrule
\end{tabular}
\vspace{1mm}
\small \textbf{为什么必须拆账户}：用户已支付未发货时退款 $\to$ 钱直接退回用户卡，\textbf{不能从商家已提现的钱里追}。"钱不出账户体系"之前商家无法 cash-out；一旦走到第三方银行追回成本极高（走法务）。
\keypoint{结算 = 把"什么时候可以给商家钱"，从拍脑袋变成\textbf{状态机}。}
\srcnote{docs/slides/04-outbox-saga.tex 事务+事件原子写；15-order-close-business.tex 关单时机}
\end{frame}

% ============================================================
\section{admin vs merchant：边界由谁守}
% ============================================================

\begin{frame}{admin 多一个动作：promote 到 merchant + 30s 缓存}
\begin{tikzpicture}[node distance=3mm and 8mm]
  \node[node-box,font=\scriptsize] (req)   {商家请求};
  \node[node-box,right=of req,node-warn,font=\scriptsize] (cache) {sync.Map\\TTL 30s};
  \node[node-box,below=of cache,font=\scriptsize] (db) {MySQL\\role='merchant'};
  \node[node-box,right=of cache,node-ok,font=\scriptsize] (ok) {放行 /merchant};
  \draw[arrow] (req) -- (cache);
  \draw[arrow,dashed] (cache) -- node[right,font=\tiny]{miss} (db);
  \draw[arrow,dashed] (db.east) to [bend right=20] node[right,font=\tiny]{回填} (cache.east);
  \draw[arrow] (cache) -- (ok);
\end{tikzpicture}
\vspace{1mm}
\begin{itemize}
  \item \textbf{现状}：\texttt{PromoteToAdmin} 硬编码 RoleAdmin。\textbf{路线}：抽象成 \texttt{PromoteToRole(target, role)}。
  \item \textbf{流程}：BD 谈妥 $\to$ 商家传资质 $\to$ admin 审核 $\to$ \texttt{promote(role="merchant")} $\to$ 解锁 \texttt{/merchant}。
  \item \textbf{对称性陷阱}（deck 14）：提权慢可接受、降权慢是真风险。封号必须显式 \texttt{InvalidateRoleCache}，否则违规商家还能操作 30s。
\end{itemize}
\srcnote{service/admin.go:47--62 PromoteToAdmin 当前写死；middleware/rbac.go:22--45,88--91 cache + Invalidate}
\end{frame}

\begin{frame}{merchant 不能跨店 = 后端必须强约束}
\begin{itemize}
  \item \textbf{前端做不到}：哪怕 UI 完全不展示别店数据，商家可以\textbf{改 URL 参数 / 改 boss\_id 字段}发请求。
  \item \textbf{鉴权做不到}：\texttt{RequireRole("merchant")} 只能保证"是某个商家"，不能保证"是这单的商家"。
  \item \textbf{必须在 dao / service 层强约束}：所有 merchant 接口的 SQL\textbf{必须}带 \texttt{WHERE boss\_id = ctx.userID}。
  \item 类比 deck 14 的安全等级：
  \begin{itemize}
    \item L1 路由前缀（区分 user/merchant/admin） —— RBAC 中间件。
    \item L2 业务约束（boss\_id 必须匹配 ctx）—— 写在 service / dao。
    \item L3 数据审计（异常聚合查询）—— 待补。
  \end{itemize}
\end{itemize}
\srcnote{repository/db/dao/product.go:67--72 已有 WHERE id AND boss\_id 模板可参考}
\end{frame}

% ============================================================
\section{压测数据告诉我们什么}
% ============================================================

\begin{frame}{商家维度查询的性能基线}
\begin{itemize}
  \item \texttt{order} 表当前 \textasciitilde 6M 行，\texttt{product} 表 \textasciitilde 956K 行（stressTest/REPORT.md L8）。
  \item C 端 \texttt{/orders/list}（WHERE user\_id 走索引）：\textbf{58{,}406 RPS / p95 5ms}。
  \item 旧版深分页（无游标）：\textbf{8.3 RPS / p95 15.95s} —— 大表上不走索引就是这个结局。
  \item \textbf{商家查"我今天的订单"}如果不建 \texttt{idx\_boss\_date(boss\_id, created\_at)} $\to$ 直接对标深分页那一行。
\end{itemize}
\keypoint{先建索引，再开接口。否则上线第一天就被自家商家 DoS。}
\srcnote{stressTest/REPORT.md 各链路压测表（行 3 vs 行 9）}
\end{frame}

% ============================================================
\section{遗留风险与路线图}
% ============================================================

\begin{frame}{当前的洞 + 修补顺序 + 边界由谁定}
\begin{columns}[T,onlytextwidth]
\column{0.55\textwidth}
\textbf{洞}：
\begin{itemize}\setlength{\itemsep}{0.5mm}
  \item \textbf{隐式商家}：任何用户调一次 product/create $\to$ 自己变商家。
  \item \textbf{无商家接口}：自助看数据走客服或 admin。
  \item \textbf{无 boss\_id 索引}：查询一上来全表扫。
  \item \textbf{无资质 / 结算表}：法务 / 财务都不闭环。
\end{itemize}
\textbf{修补顺序}（按 ROI）：
\begin{enumerate}\setlength{\itemsep}{0.5mm}
  \item RoleMerchant + RequireRole $\to$ 锁 product/create。
  \item idx\_boss\_date 索引 + 三个读接口。
  \item 库存告警 outbox + 钉钉机器人。
  \item 资质审核表 + admin 审核流程。
  \item settlement 表 + 提现接口（财务对接）。
\end{enumerate}
\column{0.42\textwidth}
\scriptsize
\textbf{边界由谁定}：
\begin{tabular}{@{}ll@{}}
\toprule
决定 & 由谁负责 \\
\midrule
谁能成为商家     & BD + 法务 \\
商家可见字段     & 产品 + 风控 \\
库存告警阈值     & 商品 + 数据 \\
结算冻结期       & 财务 + 风控 \\
提现单日限额     & 财务 + 风控 \\
封禁商家触发     & 运营 + 法务 \\
\bottomrule
\end{tabular}
\end{columns}
\vspace{1mm}
\small \textbf{最终目标}：商家自助看数据 / 上架走审核 / 库存有告警 / 结算有路径 / 出问题平台先冻结再追责。
\keypoint{商家不是"特殊的 user"，是\textbf{独立的角色}。BossID 字段在了 4 年，是时候把这件事补上。}
\end{frame}

% ============================================================
\appendix
% ============================================================

\begin{frame}{Q\&A（上）：产品 / BD 角度}
\textbf{Q1}：为什么不直接给商家开 admin 权限？\\
\hspace{4mm}\small A：admin 能看全平台数据，商家拿到 admin = 当场翻车（看到别家销售）。RBAC 三层墙的意义就是\textbf{最小权限}。

\vspace{2mm}
\textbf{Q2}：merchant 看自店订单万一查询慢怎么办？\\
\hspace{4mm}\small A：同 deck 13 思路 —— 先建 \texttt{idx\_boss\_date} 索引、再上游标分页（last\_id）；首页加 5min Redis 缓存兜底。C 端 orders/list 跑到 58K RPS，商家版数据量更小，吞吐预期类似。

\vspace{2mm}
\textbf{Q3}：商家"品牌主页"和"店铺主页"是不是一回事？\\
\hspace{4mm}\small A：业务上不是。品牌主页是 C 端引流入口（公开 GET，匿名可见），店铺主页是商家自助后台（authed + RequireRole）。数据源可同源，但路由 / 鉴权 / 缓存策略完全不同。
\end{frame}

\begin{frame}{Q\&A（下）：法务 / 风控角度}
\textbf{Q4}：第三方店铺和自营商品怎么区分？\\
\hspace{4mm}\small A：自营 = BossID 指向"平台账号"（一个特殊 user\_id）；第三方 = BossID 指向真实商家。聚合查询时一条 WHERE 就能区分。这也是为什么必须保留 BossID 而不是把自营 boss\_id 写成 0。

\vspace{2mm}
\textbf{Q5}：资质审核怎么落地？OCR 还是人工？\\
\hspace{4mm}\small A：OCR 做\textbf{字段抽取}（营业执照号 / 法人 / 经营范围），但\textbf{最终拍板必须人工} —— 营业执照能 PS、品牌授权能伪造。OCR 是节省人工录入，不是替代审核。落地：商家上传 $\to$ OCR 预填 $\to$ admin 审核员复核 $\to$ promote 到 merchant。

\vspace{2mm}
\textbf{Q6}：一个店铺多个员工账号怎么权限拆？\\
\hspace{4mm}\small A：MVP 阶段不拆，单店铺 = 单账号。等业务量足够大再加 \texttt{merchant\_employee} 表 + 子角色（运营 / 客服 / 财务）。当前不要预先抽象 —— 角色越多越容易出"谁能批准谁"的扯皮（deck 14 同样论点）。
\end{frame}

\begin{frame}{代码位置一览}
\begin{description}
  \item[BossID @ order]\srcnote{repository/db/model/order.go:12}
  \item[BossID @ product]\srcnote{repository/db/model/product.go:23}
  \item[BossID @ cart]\srcnote{repository/db/model/cart.go:10}
  \item[BossID @ favorite]\srcnote{repository/db/model/favorite.go:11--12}
  \item[BossId @ skill\_product]\srcnote{repository/db/model/skill\_product.go:6}
  \item[现有 admin group（merchant 模板）]\srcnote{routes/router.go:135--141}
  \item[product/create 当前位置]\srcnote{routes/router.go:60}
  \item[RequireRole 中间件]\srcnote{middleware/rbac.go:48--86}
  \item[InvalidateRoleCache]\srcnote{middleware/rbac.go:88--91}
  \item[PromoteToAdmin（待泛化）]\srcnote{service/admin.go:47--62}
  \item[outbox publisher（告警通道）]\srcnote{service/outbox/publisher.go:65--86}
  \item[\textbf{merchant group 待新增}]\srcnote{routes/router.go:135 后插入新 group}
\end{description}
\small 配套压测数据：stressTest/REPORT.md 链路压测表（C 端 orders/list 行 vs 深分页行）。
\end{frame}

\end{document}
